\documentclass[10pt,a4paper]{amsart}
\usepackage{amsmath,amssymb,amsthm,graphicx}
% For commutative diagrams we use Xy-pic
\usepackage[all]{xy}

% We use 2cell for 2-commutative diagrams.
\xyoption{2cell}
% Theorem environments.
%
\theoremstyle{plain}
\newtheorem{theorem}[subsection]{Theorem}
\newtheorem{proposition}[subsection]{Proposition}
\newtheorem{lemma}[subsection]{Lemma}

\theoremstyle{definition}
\newtheorem{definition}[subsection]{Definition}
\newtheorem{example}[subsection]{Example}
\newtheorem{exercise}[subsection]{Exercise}
\newtheorem{situation}[subsection]{Situation}

\theoremstyle{remark}
\newtheorem{remark}[subsection]{Remark}
\newtheorem{remarks}[subsection]{Remarks}

\numberwithin{equation}{subsection}

% Macros
%
\def\lim{\mathop{\mathrm{lim}}\nolimits}
\def\colim{\mathop{\mathrm{colim}}\nolimits}
\def\Spec{\mathop{\mathrm{Spec}}}
\def\Hom{\mathop{\mathrm{Hom}}\nolimits}
\def\Ext{\mathop{\mathrm{Ext}}\nolimits}
\def\SheafHom{\mathop{\mathcal{H}\!\mathit{om}}\nolimits}
\def\SheafExt{\mathop{\mathcal{E}\!\mathit{xt}}\nolimits}
\def\Sch{\mathit{Sch}}
\def\Mor{\mathop{\mathrm{Mor}}\nolimits}
\def\Ob{\mathop{\mathrm{Ob}}\nolimits}
\def\Sh{\mathop{\mathit{Sh}}\nolimits}
\def\NL{\mathop{N\!L}\nolimits}
\def\CH{\mathop{\mathrm{CH}}\nolimits}
\def\proetale{{pro\text{-}\acute{e}tale}}
\def\etale{{\acute{e}tale}}
\def\QCoh{\mathit{QCoh}}
\def\Ker{\mathop{\mathrm{Ker}}}
\def\Im{\mathop{\mathrm{Im}}}
\def\Coker{\mathop{\mathrm{Coker}}}
\def\Coim{\mathop{\mathrm{Coim}}}

% Boxtimes
%
\DeclareMathSymbol{\boxtimes}{\mathbin}{AMSa}{"02}

%
% Macros for moduli stacks/spaces
%
\def\QCohstack{\mathcal{QC}\!\mathit{oh}}
\def\Cohstack{\mathcal{C}\!\mathit{oh}}
\def\Spacesstack{\mathcal{S}\!\mathit{paces}}
\def\Quotfunctor{\mathrm{Quot}}
\def\Hilbfunctor{\mathrm{Hilb}}
\def\Curvesstack{\mathcal{C}\!\mathit{urves}}
\def\Polarizedstack{\mathcal{P}\!\mathit{olarized}}
\def\Complexesstack{\mathcal{C}\!\mathit{omplexes}}
% \Pic is the operator that assigns to X its picard group, usage \Pic(X)
% \Picardstack_{X/B} denotes the Picard stack of X over B
% \Picardfunctor_{X/B} denotes the Picard functor of X over B
\def\Pic{\mathop{\mathrm{Pic}}\nolimits}
\def\Picardstack{\mathcal{P}\!\mathit{ic}}
\def\Picardfunctor{\mathrm{Pic}}
\def\Deformationcategory{\mathcal{D}\!\mathit{ef}}
\usepackage{cleveref}
\crefname{originallemma}{original lemma}{original lemmas}
\Crefname{originallemma}{Original lemma}{Original lemmas}
\makeatletter
\newlabel{derived-lemma-leray-acyclicity}{{015E}{source}{}{sourcepointer-1}{}}
\newlabel{derived-lemma-leray-acyclicity@cref}{{[originallemma][0][]015E}{[1][0][]source}{}{}{}}
\makeatother
\crefname{originalequation}{original equation}{original equations}
\Crefname{originalequation}{Original equation}{Original equations}
\makeatletter
\newlabel{equation-base-change-diagram}{{02KF}{source}{}{sourcepointer-2}{}}
\newlabel{equation-base-change-diagram@cref}{{[originalequation][0][]02KF}{[1][0][]source}{}{}{}}
\makeatother
\crefname{originalreference}{original reference}{original references}
\Crefname{originalreference}{Original reference}{Original references}
\makeatletter
\newlabel{lemma-proper-base-change-f-star}{{0A3U}{source}{}{sourcepointer-3}{}}
\newlabel{lemma-proper-base-change-f-star@cref}{{[originalreference][0][]0A3U}{[1][0][]source}{}{}{}}
\makeatother
\crefname{originalreference}{original reference}{original references}
\Crefname{originalreference}{Original reference}{Original references}
\makeatletter
\newlabel{lemma-relative-colimit}{{09Z1}{source}{}{sourcepointer-4}{}}
\newlabel{lemma-relative-colimit@cref}{{[originalreference][0][]09Z1}{[1][0][]source}{}{}{}}
\makeatother
\title{Cohomology commutes with base change for a proper morphism if and only if acyclic injective test sheaves stay acyclic after base change}
\author{The Stacks Project Authors}
\date{stacks-project:etale-cohomology.tex:local-snapshot-20261011; GNU FDL 1.2 or later}
\begin{document}
\maketitle
\noindent\textit{Editorial note.} All delivered mathematical text is the author's own, copied from the pinned Stacks Project chapter file etale-cohomology.tex; no statement or proof step was rewritten, reordered or removed.\par\medskip
\noindent\textit{Editorial note.} The counted proof is the author's complete argument: the immediate direction, the reduction to sheaves annihilated by an integer via commutation with filtered colimits, the reduction to sheaves annihilated by a prime via the long exact sequence and the five lemma, and the closing injective-resolution argument using Leray's acyclicity lemma.\par\medskip
\noindent\textit{Editorial note.} The label lemma-proper-base-change-in-terms-of-injectives names the counted statement, and the proof environment immediately following it in the source is delivered verbatim as its proof.\par\medskip
\noindent\textit{Editorial note.} The counted statement is carried by the author's own lemma environment in the chapter source.\par\medskip
\noindent\textit{Editorial note.} Pre-open gate: python operations/agent-authoring/tools/precard\_gate.py --source etale-cohomology.tex --label lemma-proper-base-change-in-terms-of-injectives --start 18359 --end 18428 --local-id 3557 returned RESULT=OK (L3557 free, occupancy FREE, public interval clear). Candidate taken from lane-s4/halfB-priority.json.\par\medskip
\noindent\textit{Editorial note.} The statement and proof are copied from the pinned chapter file; the cited inputs are kept as external pointers with their canonical Stacks tags.\par\medskip
\noindent\textit{Allowed standard prerequisites (editorial).}\begin{itemize}
\item Higher direct images commute with filtered colimits. Same chapter, the lemma that the higher direct images of a quasi-compact and quasi-separated morphism commute with filtered colimits of sheaves. Supplies the commutation used to write both sides of the base change map as filtered colimits over the truncations F[n] of the coefficient sheaf, reducing to the case of a sheaf annihilated by an integer.
\item Leray's acyclicity lemma. Derived Categories, the lemma that a bounded below complex of acyclic objects computes the derived functor. Supplies the conclusion that f'\_*(g')\textasciicircum{}\{-1\}I\textasciicircum{}bullet computes Rf'\_*(g')\textasciicircum{}\{-1\}F for an injective resolution I\textasciicircum{}bullet of a sheaf annihilated by a prime.
\item Base change tested on pushforwards of injective sheaves. Same chapter, the lemma that base change for a sheaf annihilated by a prime follows from the vanishing of the higher direct images of the pullback of an injective sheaf. Supplies the final step identifying the map on the level of complexes with the base change map, completing the equivalence.
\end{itemize}\medskip
\ifdefined\hypertarget\hypertarget{sourcepointer-1}{}\fi
\noindent\textit{Source reference.} Original lemma 015E: sources/stacks-chow/tags.txt, source lines 1505--1505. Prerequisite pointer invoked at proof line(s) 18421; it names a standard fact of a sibling chapter that the counted argument uses, and every step of the counted proof built on it is retained verbatim in the delivered body.\par\medskip
\ifdefined\hypertarget\hypertarget{sourcepointer-2}{}\fi
\noindent\textit{Source reference.} Original equation 02KF: sources/stacks-etale-cohomology/pinned-original-source.tex, source lines 18313--18313. Prerequisite pointer invoked at proof line(s) 18367; it names a standard fact of the same chapter that the counted argument uses, and every step of the counted proof built on it is retained verbatim in the delivered body.\par\medskip
\ifdefined\hypertarget\hypertarget{sourcepointer-3}{}\fi
\noindent\textit{Source reference.} Original reference 0A3U: sources/stacks-etale-cohomology/pinned-original-source.tex, source lines 18282--18282. Prerequisite pointer invoked at proof line(s) 18427; it names a standard fact of the same chapter that the counted argument uses, and every step of the counted proof built on it is retained verbatim in the delivered body.\par\medskip
\ifdefined\hypertarget\hypertarget{sourcepointer-4}{}\fi
\noindent\textit{Source reference.} Original reference 09Z1: sources/stacks-etale-cohomology/pinned-original-source.tex, source lines 7535--7535. Prerequisite pointer invoked at proof line(s) 18391; it names a standard fact of the same chapter that the counted argument uses, and every step of the counted proof built on it is retained verbatim in the delivered body.\par\medskip
\begin{lemma}
\label{lemma-proper-base-change-in-terms-of-injectives}
Let $f : X \to Y$ be a proper morphism of schemes.
The following are equivalent
\begin{enumerate}
\item cohomology commutes with base change for $f$ (see above),
\item for every prime number $\ell$ and every injective
sheaf of $\mathbf{Z}/\ell\mathbf{Z}$-modules $\mathcal{I}$
on $X_\etale$ and every diagram (\ref{equation-base-change-diagram})
where $X' = Y' \times_Y X$ the sheaves
$R^qf'_*(g')^{-1}\mathcal{I}$ are zero for $q > 0$.
\end{enumerate}
\end{lemma}

\begin{proof}
It is clear that (1) implies (2). Conversely, assume (2) and let
$\mathcal{F}$ be a torsion abelian sheaf on $X_\etale$. Let $Y' \to Y$
be a morphism of schemes and let $X' = Y' \times_Y X$
with projections $g' : X' \to X$ and $f' : X' \to Y'$ as in
diagram (\ref{equation-base-change-diagram}).
We want to show the maps of sheaves
$$
g^{-1}R^qf_*\mathcal{F} \longrightarrow R^qf'_*(g')^{-1}\mathcal{F}
$$
are isomorphisms for all $q \geq 0$.

\medskip\noindent
For every $n \geq 1$, let $\mathcal{F}[n]$ be the subsheaf of sections
of $\mathcal{F}$ annihilated by $n$. Then
$\mathcal{F} = \colim \mathcal{F}[n]$.
The functors $g^{-1}$ and $(g')^{-1}$ commute with arbitrary colimits
(as left adjoints). Taking higher direct images along $f$ or $f'$
commutes with filtered colimits by Lemma \ref{lemma-relative-colimit}.
Hence we see that
$$
g^{-1}R^qf_*\mathcal{F} = \colim g^{-1}R^qf_*\mathcal{F}[n]
\quad\text{and}\quad
R^qf'_*(g')^{-1}\mathcal{F} =
\colim R^qf'_*(g')^{-1}\mathcal{F}[n]
$$
Thus it suffices to prove the result in case $\mathcal{F}$ is
annihilated by a positive integer $n$.

\medskip\noindent
If $n = \ell n'$ for some prime number $\ell$, then we obtain a short
exact sequence
$$
0 \to \mathcal{F}[\ell] \to \mathcal{F} \to
\mathcal{F}/\mathcal{F}[\ell] \to 0
$$
Observe that $\mathcal{F}/\mathcal{F}[\ell]$ is annihilated by $n'$.
Moreover, if the result holds for both $\mathcal{F}[\ell]$ and
$\mathcal{F}/\mathcal{F}[\ell]$, then the result holds by
the long exact sequence of higher direct images (and the $5$ lemma).
In this way we reduce to the case that $\mathcal{F}$ is annihilated
by a prime number $\ell$.

\medskip\noindent
Assume $\mathcal{F}$ is annihilated by a prime number $\ell$.
Choose an injective resolution $\mathcal{F} \to \mathcal{I}^\bullet$
in $D(X_\etale, \mathbf{Z}/\ell\mathbf{Z})$. Applying assumption
(2) and Leray's acyclicity lemma
(Derived Categories, Lemma \ref{derived-lemma-leray-acyclicity})
we see that
$$
f'_*(g')^{-1}\mathcal{I}^\bullet
$$
computes $Rf'_*(g')^{-1}\mathcal{F}$. We conclude by applying
Lemma \ref{lemma-proper-base-change-f-star}.
\end{proof}

\end{document}
